\section*{Alban pendant la guerre du Duc de Rohan}

C'est pour le compte du Duc de Rohan que des troupes sont recrutées
aux alentours de Saint Affrique par le vicomte de Panat. Et c'est celui
de Paulin, qui lui, se fait recruteur chez les « Albans ». Le 16 mars
1616, notre cité tombe aux mains du vicomte de Panat et, bien entendu,
est mise au pillage. Les églises ne sont pas épargnées. Celle de Miolles
est incendiée par les belligérants. Les pierres sculptées conservées en
remploi sur la façade de l'église actuelle ainsi que la clé de voûte
gardée à l'intérieur de l'édifice, pourraient bien être un souvenir de
ces temps difficiles. Apeurés, les habitants d'ici et d'à côté font
appel au gouverneur du Languedoc pour se débarrasser des pillards.
Ceux-ci finissent par déguerpir, mais contre bonne rançon et ce n'est
point là la fin.

L'accalmie sera de courte durée puisqu'en 1621, les hostilités ont
repris. C'est le moment -- en mars 1621 -- où Guillaume d'Albert, en
besoin d'argent, vend la seigneurie d'Alban à Guillaume Passemar de
Cantelauze. Sa fille, Jeanne, bénédictine au monastère de la Salvetat,
n'aura guère à se réjouir dans ses oraisons de ce que son père s'est
coiffé du nouveau titre de baron d'Alban. Les troubles, en effet,
recommencent. Les artisans du duc de Rohan s'emparent par la force de
Massuyès, Teillet, Paulin ; Alban trouve plus prudent de composer avec
la racaille et de capituler sans coup férir. C'est sage. Les finances
sont en piteux état. Les consuls d'Alban ont demandé (et obtenu) au
percepteur diocésain la remise totale des impositions.